\providecommand{\iquest}[1]{¿#1}
\section{Gestión del proyecto}
\sloppy

\subsection{Roles y responsabilidades del equipo}
La organización del trabajo sigue los principios de integración y definición del alcance del PMBOK \cite{pmbok2021}. La tabla~\ref{tab:roles} reparte las responsabilidades verificables de la entrega 1.

\begin{table}[htbp]
\centering
\footnotesize
\begin{tabularx}{\linewidth}{|p{3.4cm}|X|}
\hline
\thdr{Integrante} & \thdr{Responsabilidades en la entrega 1}\\
\hline
Robinson Espinoza Jeanpierre (l\'ider) & Visi\'on, alcance y objetivos; relaci\'on con la organizaci\'on; consentimiento y custodia LOPDP; integraci\'on de \texttt{main.tex}, compilaci\'on y PDF del SGA.\\
\hline
Naranjo Alfonzo Anthony Alexander & Elicitaci\'on y modelado: t\'ecnicas por fuente, instrumentos, contexto y matriz poder--inter\'es en fuentes editables.\\
\hline
S\'anchez Centeno Roselyn Andreina & An\'alisis y calidad documental: partes interesadas, estrategia de participaci\'on, coherencia y trazabilidad, referencias IEEE y revisi\'on del PDF.\\
\hline
\end{tabularx}
\caption{Roles y responsabilidades del equipo.}
\label{tab:roles}
\end{table}

\subsection{Cronograma hasta la semana 17}
La tabla~\ref{tab:cronograma} planifica el trabajo desde la semana 4 hasta la semana 17 del SPA 2026--2027. Las entregas posteriores a la semana 4 se enuncian como plan de proyecto y no como resultados obtenidos en esta entrega.

\begin{table}[htbp]
\centering
\footnotesize
\begin{tabularx}{\linewidth}{|p{1.7cm}|X|p{3.2cm}|}
\hline
\thdr{Semana} & \thdr{Actividad planificada} & \thdr{Producto}\\
\hline
4 & Visi\'on, alcance y bases de la entrega 1 & \texttt{main.tex}, README, ficha y alcance.\\
\hline
5 & Elicitaci\'on con AD-01, GO-01, MES-01 y COC-01 & Notas an\'onimas y contexto validado.\\
\hline
6 & Elicitaci\'on con CAJ-01, INV-01, DEL-01 y CLI-01 & Instrumentos aplicados y perspectivas revisadas.\\
\hline
7 & An\'alisis y especificaci\'on de requisitos & Lista de requisitos y trazabilidad.\\
\hline
8 & Modelado funcional y de datos & Modelos editables exportados.\\
\hline
9 & Validaci\'on de requisitos con la organizaci\'on & Registro de cambios acordados.\\
\hline
10 & Dise\~no de prototipo & Prototipo navegable.\\
\hline
11--12 & Construcci\'on del prototipo y pruebas unitarias & Versiones incrementales.\\
\hline
13 & Pruebas de integraci\'on y usabilidad & Informe de pruebas.\\
\hline
14 & Ajustes seg\'un validaci\'on & Prototipo corregido.\\
\hline
15 & Preparaci\'on de la entrega 2 & Documentaci\'on y fuentes.\\
\hline
16 & Revisi\'on final y ensayo & PDF y repositorio listos.\\
\hline
17 & Cierre del PFC & Entrega final.\\
\hline
\end{tabularx}
\caption{Cronograma del proyecto hasta la semana 17.}
\label{tab:cronograma}
\end{table}

\subsection{Plan de comunicaci\'on con la organizaci\'on}
La tabla~\ref{tab:comunicacion} precisa con qui\'en se comunica el equipo, en qu\'e momento del plan, por qu\'e canal y con qu\'e frecuencia. Las sesiones S1 a S5 corresponden al cronograma propuesto en la secci\'on~7; las fechas quedan sujetas a la disponibilidad confirmada por El Balcón Restaurant.

\begin{table}[htbp]
\centering
\footnotesize
\begin{tabularx}{\linewidth}{|p{2.4cm}|p{2.6cm}|X|p{2.7cm}|}
\hline
\thdr{Con qui\'en} & \thdr{Cu\'ando} & \thdr{Canal} & \thdr{Frecuencia}\\
\hline
AD-01 & S1 y S5 & Correo institucional \texttt{@uteq.edu.ec} y mensajer\'ia acordada & Una solicitud de agenda por sesi\'on y una validaci\'on formal al cierre de la entrega 1.\\
\hline
GO-01, MES-01, COC-01, COC-02 & S1 y S2 & Coordinaci\'on previa por AD-01; instrumentos impresos o digitales en la sesi\'on & Una convocatoria por sesi\'on y un resumen an\'onimo posterior.\\
\hline
CAJ-01, INV-01, DEL-01 & S3 & Correo institucional y cuestionario estandarizado & Una convocatoria y un recordatorio antes de la sesi\'on.\\
\hline
CLI-01, PRO-01 & S4 & Cuestionario an\'onimo distribuido por AD-01 & Una aplicaci\'on en la semana 6.\\
\hline
Equipo ISR-401 & Despu\'es de cada sesi\'on y al cierre & \texttt{evidencias/REGISTRO\_CONTACTO.md} y registro privado en el manifiesto & Un asiento por contacto y una revisi\'on final antes de la entrega.\\
\hline
\end{tabularx}
\caption{Plan de comunicaci\'on con la organizaci\'on y el equipo.}
\label{tab:comunicacion}
\end{table}

El l\'ider concentra las solicitudes y devuelve los acuerdos; los dem\'as integrantes no contactan a la organizaci\'on por separado sin coordinaci\'on. No se usan correos personales de participantes y cada contacto se anota en \texttt{evidencias/REGISTRO\_CONTACTO.md}, cuyo respaldo privado se registra en el manifiesto con SHA-256.

\subsection{Registro de riesgos}
La tabla~\ref{tab:riesgos} recoge riesgos espec\'ificos de la organizaci\'on y del contexto acad\'emico, con mitigaciones verificables.

\begin{table}[htbp]
\centering
\footnotesize
\begin{tabularx}{\linewidth}{|l|X|p{1.9cm}|X|}
\hline
\thdr{ID} & \thdr{Riesgo} & \thdr{Impacto} & \thdr{Mitigaci\'on}\\
\hline
R1 & Disponibilidad limitada de AD-01 por atenci\'on del local & Alto & Agendar con antelaci\'on y preparar res\'umenes breves para validaci\'on.\\
\hline
R2 & Ruido e interrupciones en sala durante la observaci\'on & Medio & Observar en turnos de baja demanda y registrar solo c\'odigos de rol.\\
\hline
R3 & Rotaci\'on de personal antes de la validaci\'on & Medio & Documentar el rol, no la persona, y repetir la sesi\'on cr\'itica con el reemplazo.\\
\hline
R4 & Sensibilidad de datos de clientes y proveedores & Alto & Excluir identificadores personales y aplicar anonimizaci\'on desde el instrumento.\\
\hline
R5 & Conectividad insuficiente en el local & Medio & Dise\~nar instrumentos impresos y probar la soluci\'on con datos locales.\\
\hline
R6 & Dependencia de una sola fuente para el flujo de pedidos & Medio & Contrastar entrevista, observaci\'on y cuestionario.\\
\hline
R7 & Carga acad\'emica concurrente en semana 16 & Medio & Cerrar fuentes y compilaci\'on antes de la semana 15.\\
\hline
\end{tabularx}
\caption{Registro de riesgos de la entrega 1.}
\label{tab:riesgos}
\end{table}

\subsection{\'Etica, consentimiento y custodia}
El tratamiento de datos se rige por la LOPDP \cite{lopdp2021} y por el C\'odigo de \'Etica de la ACM \cite{acm2018}. Se emplean c\'odigos de rol, se excluyen c\'edulas, tel\'efonos, correos personales, audios y firmas del repositorio p\'ublico, y los cuatro originales firmados se recibieron como fotografías el 21/11, se custodian en la carpeta institucional privada \texttt{@uteq.edu.ec} y se registran como EV-14 a EV-17 con el formato \texttt{AAAA-MM-DD\_Rol\_Codigo\_Consentimiento}, sin nombres personales. La fecha manuscrita de cada formulario no puede leerse con fiabilidad en la copia digital, por lo que no se afirma que el consentimiento haya sido previo a las sesiones. El modelo en blanco se entrega en \texttt{evidencias/MODELO\_CONSENTIMIENTO\_INFORMADO\_SGR1.tex}; el repositorio incluye el modelo en blanco, sin datos personales, y remite al manifiesto para el registro de los originales firmados. La carta de aceptaci\'on y el registro de primer contacto siguen la misma regla. Los audios y videos de las cuatro entrevistas, las transcripciones y las actas permanecen en la carpeta privada; el repositorio solo publica sus nombres de archivo, fechas, duraciones y huellas SHA-256. Este procedimiento evita exponer datos personales y permite al docente verificar la custodia sin acceder a los originales desde un repositorio abierto.

\subsection{Verificaci\'on de la evidencia de campo}
El manifiesto \texttt{evidencias/MANIFIESTO.md} registra 21 evidencias privadas (EV-01 a EV-21) con nombre de archivo, fecha, duraci\'on cuando aplica y huella SHA-256. Las evidencias EV-01 y EV-02 documentan la aceptaci\'on de la organizaci\'on y el primer contacto; EV-03 a EV-06 son las transcripciones corregidas; EV-07 a EV-09 son adendums por \'items omitidos en la transcripci\'on; EV-10 a EV-13 son los videos; EV-14 a EV-17 son los formularios de consentimiento firmados; y EV-18 a EV-21 son los audios. Ning\'un archivo identificable se sube al repositorio.

La verificaci\'on se apoya en tres criterios: (i) que el c\'odigo \texttt{ENTR-xx} sea el mismo en transcripci\'on, video, audio y formulario; (ii) que la fecha informada por el equipo y la duraci\'on del archivo coincidan con el registro; y (iii) que la huella SHA-256 permita comprobar que el archivo privado no cambi\'o. Cuando el nombre original de una grabaci\'on discrepaba con la fecha informada, se conserv\'o la fecha informada y se dej\'o constancia de la discrepancia en \texttt{INTEGRIDAD.md}. Este procedimiento no subsana la falta de consentimiento previo verificable; por eso el informe usa \'unicamente par\'afrasis an\'onimas y no publica citas textuales de las participantes.
